\section{Eigenvalues and Diagonalisation}

\subsection{Fundamental Theorem of Algebra}
\begin{theorem} (Fundamental Theorem of Algebra)

  Let $f(x) \in \mathbb{C}[x]$ be a non-constant polynomial of degree $d > 0$. The following equivalence statements hold:
  \begin{itemize}
    \item $f(x)$ has at least one complex root
    \item $f(x)$ has precisely $d$ roots counted with multiplicity
    \item $f(x)$ factorizes as a product of linear terms
  \end{itemize}
\end{theorem}

\subsection{Eigenvalues and Characteristic Polynomials of Square Matrices}

\begin{definition}
  Let $A \in M_n(\mathbb{F})$. A scalar $\lambda \in \mathbb{F}$ is an {\bf eigenvalue} of $A$ if there exists a non-zero vector $v$ such that
  \[A {\bf v} = \lambda {\bf v}.\]

In this case the vector $v$ is an {\bf eigenvector} of $A$ with {\bf eigenvalue} $\lambda$ or a {\bf $\lambda$-eigenvector}.
\end{definition}

\begin{remark}
    If $v$ is an eigenvector with eigenvector $\lambda$ then every nonzero scalar multiple $a v$ of $v$ is also a $\lambda$-eigenvector.
\end{remark}

There is a method for computing eigenvectors!

\begin{definition}
    The characteristic polynomial of a matrix $A \in M_n(\mathbb{F})$ is 
    \[\Delta_A(t) = \det (A - tI) \in \mathbb{F}[t].\]

    Where $I \in M_n(\mathbb{F})$ is the identity matrix $I = \begin{pmatrix}
    1 & ... & 0\\
    \vdots & \ddots & 0\\
    0 & ... & 1\\
    \end{pmatrix}$
\end{definition}

\begin{prop}
    The eigenvalues of $A$ are the roots of $\Delta_A(t)$ in $\mathbb{F}$.
\end{prop}
\begin{proof}
    Let $\lambda \in \mathbb{F}$. Then the following are equivalent:
    \begin{itemize}
      \item there exists non-zero $v \in \mathbb{F}^n$ such that $A v = \lambda v$.
      \item there exists non-zero $v \in N_{A - \lambda I}$ i.e. $(A -\lambda I) v = 0$
      \item $\ker \phi_{A - \lambda I} \not = 0$
      \item $\phi_{A - \lambda I}$ is not bijective
      \item $A - \lambda I$ is not invertible
      \item $\det (A - \lambda I) = 0$
      \item $\lambda$ is a solution of $\Delta_A(t) = 0$
    \end{itemize}
\end{proof}

\begin{lemma}
    If $A \in M_n(\mathbb{F})$ then $\Delta_A(t)$ is a degree $n$ polynomial, with leading coefficient $(-1)^n$. Therefore $\Delta_A(t)$ has t most $n$ roots.
    If $\mathbb{F} = \mathbb{C}$ then $\Delta_A(t)$ has precisely $n$ roots counted with multiplicity.
\end{lemma}
\begin{proof}
    \[\begin{vmatrix}
      a_{11} - t & a_{12} & ... & a_{1n}\\
      a_{21} & a_{22} - t &  & \vdots \\
    \vdots &  & \ddots & \vdots \\
    a_{n1} & ... & ... & a_{nn} - t\\
    \end{vmatrix}.\]
    has degree $n$ with the term $t^n$ arising from the summand $\sigma = Id$ permutation.

    The summand is $(a_{11} - t)(a_{22} - t)...(a_{nn} - t)$ so the leading term is $(-t)^n$, and leading coefficient $(-1)^n$.

    Hence this polynomail has maximum $n$ roots, with exactly $n$ over $\mathbb{C}$.

\end{proof}

\begin{lemma}
    Similar matrices have the same characteristic polynomial, hence the same eigenvalues.
\end{lemma}
\begin{proof}
  Suppose $A = P B P^{-1}$. Then
  \begin{align*}
    \Delta_A(t) &= \det (A - t I) \\ &= \det(PBP^{-1} - tI) \\
                &= \det(P (B - tI)P^{-1}) \\ &= \det(B - tI) = \Delta_B(t).
  .\end{align*}
\end{proof}

\begin{definition}
    Let $V$ be a vector space over $\mathbb{F}$, and $\phi : V \rightarrow V$ be a linear map.
    \begin{itemize}
      \item $\lambda \in \mathbb{F}$ is an {\bf eigenvalue} of $\phi$ if there exists nonzero $v \in V$ such that
        \[\phi(v) = \lambda v.\]
        Then $v$ is a $\lambda-${\bf eigenvector} of $\phi$.
      \item The $\lambda$-eigenspace of $\phi$ is
        \begin{align*}
          E_\phi (\lambda) &= \{ v \in V \ | \ \phi(v) = \lambda v \} \\
                           &= \ker(\phi - \lambda \id_V)
        .\end{align*}

    \end{itemize}
\end{definition}

\begin{lemma}
    Let $A \in M_n(\mathbb{F})$, $V$ be a vector space over $\mathbb{F}$, $\phi : V \rightarrow V$ be a linear map. If $A$ is the matrix of $\phi$ in some basis $\alpha$ then $A,\phi$ have the same eigenvalues and eigenvectors.
\end{lemma}

\begin{definition}
The {\bf characteristic polynomial} of $\phi : V \rightarrow V$ is
    \[\Delta_\phi(t) = \det(\phi - t \id_V).\]
    Equivalently, if $A$ represents $\phi$ in some basis $\Delta_\phi(t) := \Delta_A(t)$.
\end{definition}

\begin{remark}
    The characteristic polynomial $\phi$ is well defined because similar matrices have equal determinants.
\end{remark}

\subsection{Diagonalisation}

Let $\phi : V \rightarrow V$ be a linear map and $A \in M_n(\mathbb{F})$, the matrix of $\phi$ with respect to some basis $v_1,...,v_n$. Then the following are equivalent.
\begin{itemize}
  \item $v_i$ is a $\lambda_i-$eigenvector for $i = 1,...,n$.
  \item $A$ is diagonal and $A = \begin{pmatrix}
  \lambda_1 & \dots & 0 \\
   \vdots & \ddots & \vdots\\
   0 & \dots & \lambda_n \\ 
  \end{pmatrix}$
\end{itemize}

\begin{definition} $ $
  \begin{itemize}
  \item  A linear operator $\phi : V \rightarrow V$ is {\bf diagonalisable} if there is a basis consisting of eigenvectors of $\phi$. In this case we call it an {\bf eigenbasis}.
  \item $A \in M_n(\mathbb{F})$ is {\bf diagonalisable} (over $\mathbb{F}$) if there is an invertible $P \in M_n(\mathbb{F})$ such that
      \[P^{-1}AP \quad \text{ is diagonal}.\]
      We say that $A$ is {\bf diagonalised} by $P$.
  \end{itemize}
\end{definition}

\begin{remark}
    $A \in M_n(\mathbb{F})$ is diagonalisable if and only if $\phi_A : \mathbb{F}^n \rightarrow \mathbb{F}^n$ is diagonalisable.

\end{remark}
